\conclusion

В работе рассмотрена задача верификации соответствия трасс исполнения программы
формальной спецификации, заданной на языке TLA+. Показано, что данная задача
возникает при необходимости поддерживать согласованность между реализацией и
формальной моделью в процессе эволюции программной системы.

В ходе работы была предложена методика верификации, основанная на
инструментировании программной реализации, построении трассы наблюдаемых шагов
исполнения и её последующей интерпретации как системы ограничений на переходы
формальной модели. Такой подход позволяет сопоставлять поведение программы со
спецификацией не по исходному коду напрямую, а по наблюдаемым эффектам
выполнения, выраженным в терминах логических переменных модели.

Основным теоретическим результатом работы стала строгая математическая
постановка задачи верификации в терминах системы переходов, множества
допустимых поведений модели, множества поведений, согласованных с трассой, и
предиката согласованности шага. Показано, что рассматриваемая задача является
задачей разрешимости и сводится к проверке непустоты пересечения множеств
допустимых поведений модели и поведений, согласованных с наблюдаемой трассой.

Для решения поставленной задачи введена последовательность множеств
достижимых состояний, согласованных с префиксами трассы. Доказано, что каждое
такое множество содержит в точности те состояния модели, которые достижимы
после согласования соответствующего префикса наблюдений. На этой основе получен
критерий существования согласованного поведения и построен алгоритм проверки
согласованности трассы с моделью.

Практическая часть работы показала, каким образом полученная математическая
постановка реализуется средствами TLA+ и TLC. Рассмотрено соответствие между
переходной системой и TLA+-спецификацией, между наблюдаемой трассой и
ограниченной спецификацией, а также между рекуррентным построением множеств
достижимых состояний и обходом пространства состояний, выполняемым TLC.
Показано, что предложенная схема может быть реализована на примере
спецификации протокола двухфазной фиксации транзакции.
